\SourceLeaf{376}%
\needspace{10\baselineskip}
\begin{center}BEMERKUNGEN.\end{center}

\textsc{Gauss} hat die vorstehende Notiz auf dem Einsatzblatte eines Sammelbandes niedergeschrieben, der einige seiner Abhandlungen aus den Jahren 1799 bis 1813 enthält.

Die Integration nach $\lambda$ beruht auf der bekannten Formel:
\[
\int_0^\pi\frac{d\psi}{(u+v\cos\psi)^2}=\frac{\pi u}{(u+v)^{\frac32}(u-v)^{\frac32}};
\]
dabei ist
\[
\begin{aligned}
\psi&=2\lambda\\
u&=\tfrac12(a^2+b^2)\cos\varphi^2+c^2\sin\varphi^2\\
v&=\tfrac12(a^2+b^2)\cos\varphi^2
\end{aligned}
\]
zu setzen, was der Substitution
\[
\xi=\cos\varphi\cos\lambda,\qquad\eta=\cos\varphi\sin\lambda,\qquad\zeta=\sin\varphi
\]
entspricht.\par\nobreak\hfill\textsc{Stäckel.}\par
